\documentclass[11pt]{article}

\usepackage[margin=1in]{geometry}
\usepackage{amsmath,amssymb,amsthm}
\usepackage{graphicx}
\usepackage[hidelinks]{hyperref}

\newtheorem{theorem}{Theorem}
\newtheorem{lemma}{Lemma}
\newtheorem{corollary}{Corollary}
\newtheorem{remark}{Remark}

\newcommand{\Fix}{\operatorname{Fix}}
\newcommand{\Lip}{\operatorname{Lip}}
\newcommand{\dist}{\operatorname{d}}

\title{Uniform H\"older Nonexpansive Maps on Decomposable Banach Balls}
\author{Automatic Math Discovery Run \texttt{fa\_banach\_001}}
\date{June 12, 2026}

\begin{document}
\maketitle

\section*{Source}

Cleon S. Barroso and Valdir Ferreira,
\emph{H\"older-contractive mappings, nonlinear extension problem and fixed
point free results}, arXiv:2207.03057.

The final open questions ask whether \(B_X\) fails the fixed point property
for uniformly \(\alpha\)-H\"older nonexpansive maps with null minimal
displacement when, among other cases, \(X\) is isomorphic to a Hilbert space
or \(X\) is reflexive with an unconditional basis.

\begin{center}
\includegraphics[width=\textwidth]{figures/open_problem_crop.png}
\end{center}

\section*{Result}

\begin{theorem}
\label{thm:main}
Let \(X\) be an infinite-dimensional Banach space that splits as
\[
  X=Y\oplus Z
\]
with \(Y\) and \(Z\) closed infinite-dimensional subspaces.  Then, for every
\(\alpha\in(0,1)\), there is a map \(V:B_X\to B_X\) such that
\[
  \Fix(V)=\varnothing,\qquad
  \dist(V,B_X)=0,
\]
and
\[
  \|V^n x-V^n y\|\leq \|x-y\|^\alpha
  \qquad(x,y\in B_X,\ n\geq 1).
\]
Thus \(V\) is uniformly \(\alpha\)-H\"older nonexpansive.
\end{theorem}

\begin{corollary}
\label{cor:source-questions}
For every \(\alpha\in(0,1)\), the source paper's final open questions have
affirmative answers in the following cases:
\begin{enumerate}
\item \(X\) is isomorphic to an infinite-dimensional Hilbert space;
\item \(X\) has an unconditional basis, in particular when \(X\) is reflexive
and has an unconditional basis.
\end{enumerate}
\end{corollary}

\section*{Input Lemma From The Run}

We use the decomposable-space packet already established in this run:
if \(X=Y\oplus Z\) with both summands infinite-dimensional, then there is a
map \(U:B_X\to B_X\) such that
\[
  \Fix(U)=\varnothing,\qquad
  \dist(U,B_X)=0,\qquad
  M:=\sup_{n\geq 1}\Lip(U^n)<\infty .
\]
This is the result recorded in
\[
\texttt{solutions/partial/2307.12958\_decomposable\_tube\_uniformly\_lipschitz\_q2/}.
\]
The present packet only adds the H\"older-nonexpansive scaling step.

\section*{Scaling Lemma}

\begin{lemma}
\label{lem:scaling}
Let \(X\) be a Banach space.  Suppose \(U:B_X\to B_X\) satisfies
\[
  \Fix(U)=\varnothing,\qquad
  \dist(U,B_X)=0,\qquad
  M:=\sup_{n\geq 1}\Lip(U^n)<\infty .
\]
Then for every \(\alpha\in(0,1)\) there is a uniformly
\(\alpha\)-H\"older nonexpansive map \(V:B_X\to B_X\) with
\(\Fix(V)=\varnothing\) and \(\dist(V,B_X)=0\).
\end{lemma}

\begin{proof}
Let \(C=2M\) and put
\[
  \eta=\min\{2,\, C^{-1/(1-\alpha)}\}.
\]
Choose \(r>0\) so small that \(2r\leq \eta^\alpha\).  In particular
\(r<1\).  Let \(R_r:B_X\to rB_X\) be the radial retraction
\[
  R_r(x)=
  \begin{cases}
  x,& \|x\|\leq r,\\
  r\,x/\|x\|,& \|x\|>r.
  \end{cases}
\]
This retraction is \(2\)-Lipschitz on every normed space.

Define the scaled map \(S:rB_X\to rB_X\) by
\[
  S(z)=r\,U(z/r).
\]
Then \(S^n(z)=r\,U^n(z/r)\), so \(\Lip(S^n)\leq M\) for every \(n\geq1\).
Moreover, \(S\) is fixed-point free and has null minimal displacement on
\(rB_X\).

Now set
\[
  V=S R_r:B_X\to B_X.
\]
Since \(S(rB_X)\subseteq rB_X\) and \(R_r\) is the identity on \(rB_X\),
\[
  V^n=S^n R_r\qquad(n\geq1).
\]
For \(x,y\in B_X\), write \(d=\|x-y\|\).  Then
\[
  \|V^n x-V^n y\|
  \leq \min\{M\|R_r x-R_r y\|,\,2r\}
  \leq \min\{Cd,\,2r\}.
\]
If \(d\leq \eta\), then \(Cd\leq d^\alpha\) by the definition of \(\eta\).
If \(d\geq\eta\), then \(2r\leq\eta^\alpha\leq d^\alpha\).  Hence
\[
  \|V^n x-V^n y\|\leq d^\alpha
\]
for all \(x,y\in B_X\) and every \(n\geq1\).

There is no fixed point: if \(Vx=x\), then \(x\in rB_X\), so \(R_r x=x\)
and \(Sx=x\), contradicting \(\Fix(S)=\varnothing\).

Finally, since \(S\) has null minimal displacement on \(rB_X\), there are
\(z_j\in rB_X\) such that \(\|S z_j-z_j\|\to0\).  For these points
\(R_r z_j=z_j\), so \(\|V z_j-z_j\|\to0\).  Thus
\(\dist(V,B_X)=0\).
\end{proof}

\begin{proof}[Proof of Theorem~\ref{thm:main}]
Apply the input lemma from the run to obtain \(U\), then apply
Lemma~\ref{lem:scaling}.
\end{proof}

\begin{proof}[Proof of Corollary~\ref{cor:source-questions}]
If \(X\) is isomorphic to an infinite-dimensional Hilbert space, split the
Hilbert space into two closed infinite-dimensional orthogonal summands and
transport that decomposition to \(X\) through an isomorphism.  The transported
summands are closed and complemented, so Theorem~\ref{thm:main} applies.

If \(X\) has an unconditional basis, split the basis into its odd and even
subsequences.  The associated coordinate subspaces are closed,
infinite-dimensional, and complemented by the unconditional coordinate
projections.  Again Theorem~\ref{thm:main} applies.  The reflexive
unconditional-basis case from the source paper is therefore included.
\end{proof}

\section*{Limitations}

This does not answer the source paper's arbitrary reflexive case, nor the
James-space case.  The argument relies on the existence of a decomposable
two-infinite-summand structure, because that is what supplies the
fixed-point-free uniformly Lipschitzian dynamics with null minimal
displacement.

\section*{Verification Notes}

The only scalar estimate in the new argument is
\[
  \min\{Cd,2r\}\leq d^\alpha\qquad(0\leq d\leq2),
\]
with \(\eta=\min\{2,C^{-1/(1-\alpha)}\}\) and \(2r\leq\eta^\alpha\).
The accompanying script \texttt{code/verify\_scaling\_lemma.py} samples this
inequality for representative values of \(C\), \(r\), and \(\alpha\).

\end{document}
